\documentclass[12pt]{article}
\usepackage{sbc-template}
\usepackage{graphicx,url}
\usepackage[utf8]{inputenc}
\usepackage[brazil]{babel}
\usepackage{booktabs}
\usepackage{amsmath}
\usepackage{amssymb}
\usepackage{multirow}
\usepackage{array}
\usepackage{float}
\usepackage{tikz}
\usetikzlibrary{shapes.geometric,arrows.meta,positioning,fit,backgrounds,calc}
\usepackage{pgfplots}
\pgfplotsset{compat=1.18}
\usepackage{xcolor}
\definecolor{azulRI}{RGB}{0,70,127}
\definecolor{cinzaRI}{RGB}{100,100,100}
\definecolor{verdeRI}{RGB}{0,120,60}

\sloppy

\title{BirdSearch: Recuperação de Informação Acústica para\\ Identificação
de Espécies de Aves com Fusão de Evidências e Reranking}

\author{Pedro Mendes\inst{1}, Alcir Heber\inst{1}}

\address{Disciplina ICC222 -- Tópicos em Recuperação de Informação\\
  Instituto de Computação -- Universidade Federal do Amazonas (UFAM)
  \email{\{pedro.mendes, alcir.heber\}@icomp.ufam.edu.br}
}

\begin{document}
\maketitle

% ─── Abstract (English) ───────────────────────────────────────────────────────
\begin{abstract}
This paper presents \textbf{BirdSearch}, a modular acoustic information
retrieval system for bird species identification, evaluated on the BirdCLEF
2026 benchmark (34,144 recordings, 571 species). The system investigates two
indexing strategies --- segment-based late fusion using an HNSW graph index
(E1) and early fusion via max-pooled super-embeddings with a flat index (E2)
--- combined with 14 reranking families across two neural embedding backbones:
Google Perch v2 and BirdNET v3.0 (both 1280-dimensional). On the full dataset,
E2 achieves MAP\,=\,0.8227 and P@1\,=\,0.7729 at 6.3\,ms average latency,
while E1 reaches MAP\,=\,0.8191 with better list diversity (P@5\,=\,0.2518)
at 22.6\,ms. Score-based rerankers (attention and softmax) consistently
outperform rank-fusion alternatives (RRF, Borda), with the best individual
result being MAP\,=\,0.8944 using BirdNET + attention in E1. Hybrid rankers
combining score smoothing and taxonomic reinforcement achieve comparable MAP
at sub-millisecond latency (0.3--1.0\,ms). Overlapping segmentation (50\%
hop) proves essential: its removal degrades MAP by 2--3 points across both
strategies.
\end{abstract}

% ─── Resumo (Português) ───────────────────────────────────────────────────────
\begin{resumo}
Este artigo apresenta o \textbf{BirdSearch}, um sistema modular de recuperação
de informação acústica para identificação de espécies de aves, avaliado sobre
o benchmark BirdCLEF 2026 (34.144 gravações, 571 espécies). O sistema investiga
duas estratégias de indexação --- fusão tardia por segmentos com índice HNSW
(E1) e fusão precoce via super-embeddings com max-pooling e índice plano (E2)
--- combinadas com 14 famílias de reranking sobre dois \textit{backbones}
neurais: Google Perch v2 e BirdNET v3.0 (ambos com 1280 dimensões). No dataset
completo, E2 atinge MAP\,=\,0,8227 e P@1\,=\,0,7729 com latência média de
6,3\,ms, enquanto E1 alcança MAP\,=\,0,8191 com maior diversidade de lista
(P@5\,=\,0,2518) a 22,6\,ms. Rerankers baseados em score (\texttt{attention}
e \texttt{softmax}) superam consistentemente alternativas de fusão por rank
(RRF, Borda), com o melhor resultado individual sendo MAP\,=\,0,8944 com
BirdNET\,+\,attention em E1. Rankers híbridos que combinam suavização de score
e reforço taxonômico atingem MAP comparável com latência sub-milissegundo
(0,3--1,0\,ms). A segmentação com sobreposição de 50\% mostra-se essencial:
sua remoção degrada o MAP em 2--3 pontos em ambas as estratégias.
\end{resumo}

% ══════════════════════════════════════════════════════════════════════════════
\section{Introdução}
\label{sec:intro}
% ══════════════════════════════════════════════════════════════════════════════

A identificação automática de espécies de aves a partir de gravações acústicas
é um problema de recuperação de informação (RI) multimídia com aplicações
diretas em monitoramento ambiental, conservação da biodiversidade e estudos
ecológicos de larga escala~\cite{stowell:22}. A plataforma Xeno-canto já reúne
mais de 700.000 gravações de mais de 10.000 espécies~\cite{xeno:14}, tornando
a busca eficiente por similaridade acústica um desafio de RI não trivial.

O desafio \textbf{BirdCLEF 2026} fornece um benchmark padronizado com 34.144
gravações de campo de 571 espécies, predominantemente da região Brasileira.
Dado um áudio de consulta, o objetivo é recuperar e ranquear as espécies mais
prováveis no corpus. Esse cenário distingue-se da classificação clássica por
três características: (a) o espaço de classes é fechado mas vasto (571
espécies); (b) cada gravação pode conter vocalizações de múltiplas espécies;
e (c) a duração dos arquivos é variável, exigindo estratégias de segmentação
e fusão de evidências.

Abordagens recentes de recuperação acústica de aves exploram embeddings de
redes neurais pré-treinadas~\cite{kahl:21, merriënboer:25}. No entanto, a
literatura carece de análises sistemáticas sobre: (i) como estruturar o índice
quando cada gravação gera múltiplos embeddings de segmento; e (ii) quais
estratégias de reranking melhor exploram as propriedades dos embeddings gerados
por modelos especializados em aves.

O \textbf{BirdSearch} aborda exatamente essas duas questões. Avaliamos o
impacto de duas estratégias de indexação e fusão --- \textit{late fusion} por
segmentos (E1) e \textit{early fusion} via super-embedding (E2) --- combinadas
com 14 rerankers de cinco famílias distintas, sobre dois backbones de embedding.
Os experimentos abrangem 6.835 consultas sobre um corpus de 27.309 arquivos de
treino, totalizando 363.917 segmentos no índice E1.

As principais contribuições deste trabalho são:
\begin{itemize}
  \item Uma análise sistemática do impacto do
  pré-processamento espectral e das estratégias de fusão sobre MAP, MRR,
  nDCG, P@k, R@k e métricas taxonômicas.
  \item Uma avaliação de 14 rerankers sobre dois backbones, identificando os
  melhores candidatos para diferentes critérios (precisão, cobertura,
  proximidade biológica, latência).
\end{itemize}

Este artigo está organizado da seguinte forma: a Seção~\ref{sec:ref} apresenta
o referencial teórico; a Seção~\ref{sec:fund} detalha as métricas e o
embasamento matemático; a Seção~\ref{sec:metod} descreve a metodologia e a
arquitetura; a Seção~\ref{sec:res} apresenta e discute os resultados; e a
Seção~\ref{sec:conc} conclui o trabalho.

% ══════════════════════════════════════════════════════════════════════════════
\section{Referencial Teórico}
\label{sec:ref}
% ══════════════════════════════════════════════════════════════════════════════

\subsection{Recuperação de Informação por Similaridade}

A recuperação de informação por similaridade vetorial (\textit{nearest-neighbor
retrieval}) é o paradigma central de sistemas de busca em espaços contínuos de
alta dimensão~\cite{baeza:99}. Dada uma coleção de $N$ vetores de documentos
$\{\mathbf{d}_i\} \subset \mathbb{R}^d$ e uma consulta $\mathbf{q} \in
\mathbb{R}^d$, o objetivo é encontrar os $k$ vetores mais similares a
$\mathbf{q}$ segundo uma métrica de distância ou similaridade. Em domínios de
áudio, a similaridade cosseno é a métrica predominante, pois é invariante à
magnitude do sinal e captura ângulos no espaço de características~\cite{kong:20}.

Para coleções de grande escala, a busca exata ($O(Nd)$ por consulta) torna-se
proibitiva. Índices de vizinhança aproximada (\textit{Approximate Nearest
Neighbor}, ANN) como o FAISS~\cite{johnson:21} oferecem estruturas eficientes
que reduzem a complexidade de busca mantendo alta taxa de recall. Entre eles,
o HNSW (\textit{Hierarchical Navigable Small World})~\cite{malkov:20}
destaca-se por combinar construção eficiente com latência de busca sub-linear
e recall superior a 99\% em configurações típicas.

\subsection{Embeddings Neurais para Áudio de Aves}

A extração de embeddings de áudio via redes neurais pré-treinadas transformou
o campo da bioacústica computacional. Os PANNs (\textit{Pre-trained Audio
Neural Networks})~\cite{kong:20} são redes convolucionais treinadas sobre o
AudioSet~\cite{gemmeke:17} para reconhecimento geral de eventos sonoros. Embora
versáteis, seu desempenho em domínios específicos como vocalizações de aves é
limitado pela distribuição generalista dos dados de treinamento.

O \textbf{Google Perch v2}~\cite{merriënboer:25} representa uma evolução
especializada: treinado exclusivamente sobre mais de 400.000 gravações da
plataforma Xeno-canto, o modelo combina uma arquitetura baseada em EfficientNet
com mecanismos de atenção, produzindo vetores de 1280 dimensões normalizados em
L2. Por operar em domínio específico, o Perch v2 captura padrões vocais finos
que modelos generalistas perdem.

O \textbf{BirdNET v3.0}~\cite{kahl:21}, desenvolvido pelo Cornell Lab of
Ornithology, é outra referência de embedding especializado em aves. Treinado
para classificação de 6.000+ espécies, também produz vetores de 1280 dimensões,
mas com janela nativa de 3~s a 32~kHz. A comparação entre os dois backbones
permite isolar o impacto da distribuição de treinamento sobre as estratégias
de RI.

\subsection{Fusão de Evidências em RI Multimídia}

Quando uma consulta produz múltiplas representações (e.g., segmentos de um
áudio), existem duas abordagens clássicas de fusão~\cite{baeza:99}:

\textit{Early fusion} (fusão precoce) combina as representações \textbf{antes}
da busca, produzindo uma consulta única. É eficiente e compacto, mas pode
perder sinais localizados em segmentos específicos.

\textit{Late fusion} (fusão tardia) emite múltiplas buscas independentes e
combina as \textbf{listas de resultados}. Preserva a granularidade das evidências
locais, às custas de maior latência e complexidade de combinação.

\subsection{Fusão de Rankings}

A combinação de múltiplas listas de resultados é estudada extensivamente em RI.
O \textit{Reciprocal Rank Fusion} (RRF)~\cite{cormack:09} é um método robusto
e sem parâmetros que pondera cada documento pelo inverso de sua posição em cada
lista, sem requerer normalização de scores. Formalmente:
\begin{equation}
  \text{RRF}(d, L_1,\ldots,L_n) = \sum_{i=1}^{n}
  \frac{1}{k + \text{rank}_i(d)}, \quad k = 60
  \label{eq:rrf}
\end{equation}

Estudos empíricos mostram que o RRF frequentemente supera métodos mais
sofisticados de combinação de listas~\cite{cormack:09}, sendo particularmente
robusto quando os scores brutos das listas têm distribuições heterogêneas.

\subsection{Reranking em Sistemas de RI}

O reranking é uma etapa pós-recuperação que reordena os candidatos do kNN para
maximizar algum critério de qualidade. Sistemas de busca modernos combinam
frequentemente recuperação vetorial eficiente com modelos de reranking mais
expressivos~\cite{baeza:99}. No contexto de aves, o reranking oferece uma
oportunidade única: incorporar conhecimento taxonômico (gênero, família) para
promover candidatos biologicamente plausíveis, mesmo quando a espécie exata
não é identificada com certeza.

\subsection{Métricas de Avaliação em RI}

As métricas canônicas de avaliação de sistemas de RI estão bem estabelecidas
na literatura~\cite{baeza:99, jarvelin:02}. A \textit{Average Precision} (AP)
combina precisão e recall em uma única pontuação que penaliza resultados
relevantes em posições tardias. O nDCG~\cite{jarvelin:02} introduz um desconto
logarítmico por posição, sendo mais sensível à qualidade do ranking nas posições
superiores. O MRR foca exclusivamente na posição do primeiro resultado correto.

% ══════════════════════════════════════════════════════════════════════════════
\section{Fundamentação Teórica}
\label{sec:fund}
% ══════════════════════════════════════════════════════════════════════════════

Só dizer que a gente usar as métrica de MAP,.....

%\subsection{Similaridade Cosseno e Produto Interno}

% O BirdSearch adota a similaridade cosseno como métrica de busca:
% \begin{equation}
%   \sigma_{\cos}(\mathbf{q}, \mathbf{d}) =
%   \frac{\mathbf{q} \cdot \mathbf{d}}{\|\mathbf{q}\|_2\,\|\mathbf{d}\|_2}
%   \label{eq:cosine}
% \end{equation}

% Como o Perch v2 e o BirdNET normalizam seus vetores em L2
% ($\|\mathbf{e}\|_2 = 1$), a Equação~\eqref{eq:cosine} reduz-se ao produto
% interno $\sigma = \mathbf{q} \cdot \mathbf{d}$, permitindo o uso de
% \texttt{IndexFlatIP} do FAISS com máxima eficiência.

% \subsection{Average Precision e MAP}

% Para uma consulta $q$ com rótulo $\ell$, seja $r_1,\ldots,r_k$ a lista
% recuperada e $\text{rel}(r_i) = \mathbf{1}[\text{label}(r_i) = \ell]$ a
% relevância binária. A Average Precision é:
% \begin{equation}
%   \text{AP}(q) =
%   \frac{1}{\min(|\mathcal{R}_\ell|,\,k)}
%   \sum_{i=1}^{k} \text{rel}(r_i)
%   \cdot \frac{\sum_{j \leq i} \text{rel}(r_j)}{i}
%   \label{eq:ap}
% \end{equation}
% onde $\mathcal{R}_\ell$ é o conjunto de documentos relevantes do corpus. O
% MAP é a média das APs sobre todas as consultas $Q$:
% \begin{equation}
%   \text{MAP} = \frac{1}{|Q|}\sum_{q \in Q} \text{AP}(q)
%   \label{eq:map}
% \end{equation}

% \subsection{Mean Reciprocal Rank}

% O MRR foca na posição do primeiro resultado correto:
% \begin{equation}
%   \text{MRR} = \frac{1}{|Q|}\sum_{q \in Q} \frac{1}{\text{rank}_q^*}
%   \label{eq:mrr}
% \end{equation}
% onde $\text{rank}_q^* = \infty$ (contribuição zero) se nenhum resultado
% relevante for encontrado.

% \subsection{Normalized Discounted Cumulative Gain}

% O nDCG~\cite{jarvelin:02} pondera cada posição com um desconto logarítmico:
% \begin{equation}
%   \text{DCG}@k = \sum_{i=1}^{k}
%   \frac{2^{\,\text{rel}(r_i)} - 1}{\log_2(i + 1)}
%   \label{eq:dcg}
% \end{equation}
% \begin{equation}
%   \text{nDCG}@k = \frac{\text{DCG}@k}{\text{IDCG}@k}
%   \label{eq:ndcg}
% \end{equation}
% onde IDCG@k é o DCG do ranking ideal. O desconto $1/\log_2(i+1)$ vale 1,0 na
% posição 1, cai a $\approx$\,0,63 na posição 2 e $\approx$\,0,29 na posição 10.

% \subsection{Precision@k e Recall@k}

% \begin{equation}
%   \text{P@}k = \frac{|\text{top-}k \cap \mathcal{R}_\ell|}{k}
%   \qquad
%   \text{R@}k = \frac{|\text{top-}k \cap \mathcal{R}_\ell|}
%                     {\max(1,|\mathcal{R}_\ell|)}
%   \label{eq:pr}
% \end{equation}

% O denominador de R@k depende da representação de documentos: na representação
% \textit{audio} (um vetor por arquivo), $|\mathcal{R}_\ell| = 1$ e R@1\,=\,P@1;
% na representação \textit{segment} (um vetor por janela de 5~s), $|\mathcal{R}_\ell|$
% pode chegar a centenas, tornando o recall numericamente muito menor.

\subsection{Métricas Taxonômicas}
\label{sub:tax_fund}

As métricas clássicas de RI (MAP, MRR, nDCG) tratam a relevância de forma
binária: a espécie recuperada é a correta ou não é. Para identificação de
aves, porém, \textbf{errar o gênero é muito pior do que acertar o gênero mas
errar a espécie}. Recuperar \textit{Tangara cyanicollis} quando a consulta
era \textit{Tangara chilensis} (mesma família, mesmo gênero, canto similar)
tem muito mais valor prático do que recuperar \textit{Cochlearius cochlearius}
(ave de outra ordem), mesmo que o score de similaridade seja igual.

Para capturar essa utilidade biológica, o BirdSearch define duas métricas
complementares que avaliam o \emph{quão próximos} os resultados recuperados
estão da espécie correta na árvore taxonômica:

\paragraph{genus@k.} Mede a fração dos top-$k$ resultados que pertencem ao
\textbf{mesmo gênero} que a espécie consultada:
\begin{equation}
  \text{genus@}k = \frac{1}{|Q|}\sum_{q\in Q}
  \frac{|\{r \in \text{top-}k : \text{genus}(r) = \text{genus}(q)\}|}{k}
  \label{eq:genus}
\end{equation}
Por exemplo, genus@5\,=\,0,92 significa que, em média, 4,6 dos 5 resultados
mais bem ranqueados pertencem ao mesmo gênero da espécie da consulta.

\paragraph{common@k.} Mede a fração dos top-$k$ que compartilham ao menos
uma palavra com o \textbf{nome comum} da espécie consultada:
\begin{equation}
  \text{common@}k = \frac{1}{|Q|}\sum_{q\in Q}
  \frac{|\{r \in \text{top-}k : \text{tokens}(r) \cap \text{tokens}(q) \neq \emptyset\}|}{k}
  \label{eq:common}
\end{equation}
onde \text{tokens}$(s)$ são as palavras do nome comum da espécie $s$ em inglês.
Por exemplo, se a consulta é ``Blue-naped Chlorophonia'', resultados com os
nomes ``Yellow-collared Chlorophonia'' ou ``Blue-crowned Chlorophonia'' pontuam,
pois compartilham a palavra ``Chlorophonia''. Isso captura proximidade semântica
no vocabulário ornitológico, que é frequentemente indicativo de parentesco.

Ambas as métricas valem entre 0 e 1 e são interpretadas como precisão
biologicamente relaxada: um sistema pode ter MAP\,=\,0,85 (acerta a espécie
exata em 85\% dos casos) e ainda assim ter genus@5\,=\,0,93 (em 93\% dos 5
primeiros resultados, o gênero está correto). Esse valor é relevante para
cenários de apoio a ornitólogos, onde uma lista de espécies do mesmo gênero
como sugestão já é suficiente para orientar o trabalho de campo.

\subsection{Agregadores de Embeddings}

Quando $N$ embeddings de segmento $\mathbf{e}_1,\ldots,\mathbf{e}_N \in
\mathbb{R}^{1280}$ precisam ser fundidos em um único vetor, o sistema aplica
um dos seguintes operadores (todos seguidos de normalização L2):
\begin{align}
  \mathbf{v}_{\text{mean}}   &= \ell_2\!\left(\tfrac{1}{N}\textstyle\sum_i \mathbf{e}_i\right)
  \label{eq:mean}\\
  \mathbf{v}_{\text{max}}    &= \ell_2\!\left(\max_i \mathbf{e}_i\right)
  \label{eq:max}\\
  \mathbf{v}_{\text{att}} &= \ell_2\!\left(\sum_i \alpha_i \mathbf{e}_i\right),
  \quad
  \alpha_i = \text{softmax}\!\left(\frac{\mathbf{e}_i \cdot \mathbf{c}}{\tau}\right),
  \quad \mathbf{c} = \ell_2\!\left(\tfrac{1}{N}\textstyle\sum_i\mathbf{e}_i\right)
  \label{eq:att}
\end{align}
onde $\ell_2(\mathbf{v}) = \mathbf{v}/\|\mathbf{v}\|_2$ e $\tau = 0{,}1$ é
a temperatura de atenção. O operador \textit{max} preserva o valor máximo por
dimensão, concentrando as evidências mais fortes; o \textit{attention} pondera
os segmentos por alinhamento com o centróide.

% ══════════════════════════════════════════════════════════════════════════════
\section{Metodologia}
\label{sec:metod}
% ══════════════════════════════════════════════════════════════════════════════

% \subsection{Visão Geral da Arquitetura}

% O BirdSearch é um \textit{framework} modular orientado a configuração no qual
% cada estágio do pipeline é instanciado a partir de um arquivo YAML, sem
% modificação de código. A Figura~\ref{fig:pipeline} ilustra o fluxo completo.

% \begin{figure}[ht]
% \centering
% \begin{tikzpicture}[
%     stage/.style={rectangle, rounded corners=3pt,
%                   draw=azulRI!70, fill=azulRI!8,
%                   minimum width=2.3cm, minimum height=0.75cm,
%                   font=\scriptsize\bfseries, text centered, inner sep=3pt},
%     arr/.style={-Stealth, thick, azulRI!70},
%     lbl/.style={font=\tiny\itshape, color=cinzaRI},
%     grp/.style={draw=gray!35, dashed, rounded corners=3pt, inner sep=5pt},
%     node distance=0.45cm and 0.6cm
%   ]
%   % Top
%   \node[stage] (load) {Dados};
%   \node[stage, below=of load] (split) {Split 80/20};

%   % Left branch (train/index)
%   \node[stage, below left=0.7cm and 1.1cm of split] (pt) {Pré-proc.};
%   \node[stage, below=of pt] (st) {Segmentação};
%   \node[stage, below=of st] (et) {Embedding};
%   \node[stage, below=of et] (ds) {Doc Store};
%   \node[stage, below=of ds] (ix) {Índice};

%   % Right branch (test/query)
%   \node[stage, below right=0.7cm and 1.1cm of split] (pq) {Pré-proc.};
%   \node[stage, below=of pq] (sq) {Segmentação};
%   \node[stage, below=of sq] (eq) {Embedding};
%   \node[stage, below=of eq] (fu) {Fusão};

%   % Retrieval bottom
%   \node[stage, below=1.5cm of ds] (ret) {Busca kNN};
%   \node[stage, below=of ret] (rk)  {Reranking};
%   \node[stage, below=of rk]  (ev)  {Avaliação};

%   % Arrows train
%   \draw[arr] (load) -- (split);
%   \draw[arr] (split) -| (pt);
%   \draw[arr] (pt)--(st)--(et)--(ds)--(ix);
%   % Arrows query
%   \draw[arr] (split) -| (pq);
%   \draw[arr] (pq)--(sq)--(eq)--(fu);
%   % Arrows retrieval
%   \draw[arr] (ix) |- (ret);
%   \draw[arr] (fu) |- (ret);
%   \draw[arr] (ret)--(rk)--(ev);

%   % Background groups
%   \begin{scope}[on background layer]
%     \node[grp, fit=(pt)(st)(et)(ds)(ix)] {};
%     \node[grp, fit=(pq)(sq)(eq)(fu)] {};
%   \end{scope}

%   % Labels
%   \node[lbl, left=0.05cm of pt]  {Treino};
%   \node[lbl, right=0.05cm of pq] {Teste};
% \end{tikzpicture}
% \caption{Pipeline do BirdSearch. Lado esquerdo: construção do índice a partir
% do conjunto de treino. Lado direito: processamento das consultas de teste.
% A busca kNN combina as duas ramificações.}
% \label{fig:pipeline}
% \end{figure}

% A modularidade é garantida pelo \textbf{padrão de registro}: cada estágio
% mantém uma instância de \texttt{Registry} que mapeia chaves de \textit{string}
% a classes Python (registradas via decorador \texttt{@REGISTRY.register(``nome'')}).
% O \texttt{PipelineBuilder} lê o YAML e instancia os componentes sem nenhum
% \texttt{if/elif} explícito. Substituir o indexador de HNSW para Flat, por
% exemplo, requer apenas alterar \texttt{indexing.type: hnsw} $\to$
% \texttt{indexing.type: flat} no arquivo de configuração.

\subsection{Pré-processamento}

O pré-processamento opera sobre o sinal de áudio bruto antes da segmentação.
Três módulos estão disponíveis:

% \textbf{Identidade} (\texttt{none}): passa o sinal sem modificação.

\textbf{Gating Espectral} (\texttt{spectral\_gating}): estima o piso de
ruído a partir dos 10\% de frames STFT mais silenciosos e atenua componentes
abaixo de um limiar adaptativo:
\begin{equation}
  \hat{M}(f,t) = \max\!\left(M(f,t) - \alpha \cdot N(f),\; 0\right)
  \label{eq:sg}
\end{equation}
onde $M(f,t)$ é a magnitude STFT, $N(f)$ é o piso estimado e
$\alpha=1{,}0$ é a intensidade de supressão. O sinal limpo é reconstruído
por ISTFT com \textit{overlap-add}.

\textbf{Filtro Passa-Banda} (\texttt{bandpass}): Butterworth de 4ª ordem
entre 250 e 12.000~Hz via \texttt{scipy.signal.sosfiltfilt}, cobrindo a
faixa típica de vocalizações de aves.

\subsection{Segmentação Temporal}

O Perch v2 e o BirdNet V3 opera nativamente sobre janelas de 5~s a 32~kHz (160.000 amostras).
O segmentador mais relevante neste trabalho é o de \textbf{sobreposição}
(\texttt{overlapping}): janelas deslizantes de $W=5$~s com passo $H=2{,}5$~s
(50\% de sobreposição). Para um áudio de duração $D$, o número de segmentos é:
\begin{equation}
  N_{\text{seg}} = \left\lfloor \frac{D - W}{H} \right\rfloor + 1
  \label{eq:nseg}
\end{equation}
A sobreposição garante que eventos acústicos que cruzam fronteiras de janela
sejam capturados integralmente por pelo menos um segmento. Segmentos mais
curtos que 5~s são preenchidos com zeros; os mais longos são recortados ao centro.

% \subsection{Extração de Embeddings}

% \subsubsection{Google Perch v2}
% Treinado sobre $>400.000$ gravações da Xeno-canto, o Perch v2 produz vetores
% de 1280 dimensões L2-normalizados para cada janela de 5~s a 32~kHz.
% O modelo é carregado via \texttt{kagglehub} (preferencial) ou TF-Hub. Uma
% variante \texttt{perch\_v2\_torch} exporta o SavedModel para ONNX e executa
% via PyTorch, com \textit{fallback} automático ao TF em caso de operações
% não suportadas.

% \subsubsection{BirdNET v3.0}
% Desenvolvido pelo Cornell Lab of Ornithology, o BirdNET~\cite{kahl:21} é
% treinado para classificar $>6.000$ espécies e produz vetores de 1280 dimensões.
% Sua janela nativa é de 3~s a 32~kHz; no pipeline, as janelas de 5~s são
% adaptadas por recorte ou preenchimento.

% \textbf{Cache de embeddings}: o sistema armazena os vetores computados em
% arquivos \texttt{.npz}, com chave de cache baseada em
% $(\text{path}, \text{mtime}, \text{embedder}, \text{preproc\_id}, \text{segmenter\_id})$.
% Consultas com ruído injetado sempre ignoram o cache.

\subsection{Representações de Documentos e Índices}
\label{sub:repr}

A Tabela~\ref{tab:repr} apresenta as representações de documentos disponíveis.

\begin{table}[ht]
\centering
\caption{Representações de documentos no BirdSearch.}
\label{tab:repr}
\begin{tabular}{lccc}
\toprule
\textbf{Tipo} & \textbf{Docs/arquivo} & \textbf{Tamanho (corpus)} & \textbf{Fusão indicada} \\
\midrule
\texttt{segment}   & $N_{\text{seg}}$ & 363.917 (E1) & Tardia  \\
\texttt{audio}     & 1                & 27.309 (E2)  & Precoce \\
\bottomrule
\end{tabular}
\end{table}

O índice E1 usa \textbf{HNSW}~\cite{malkov:20} com $M=32$,
\texttt{ef\_construction}\,=\,200 e \texttt{ef\_search}\,=\,64, atingindo
recall\,$>$\,0,99 e latência sub-linear ($O(\log N)$). O índice E2 usa busca
exata (\texttt{IndexFlatIP} do FAISS~\cite{johnson:21}), viável dado o tamanho
reduzido (27.309 vetores).

\subsection{Estratégias de Fusão de Evidências}
\label{sub:fusao}

Uma gravação de campo raramente tem duração exata de 5~s. Após a segmentação
com sobreposição, cada áudio gera entre 4 e 12 janelas de 5~s, cada uma com
seu próprio embedding de 1280 dimensões. O problema central é: \emph{como
usar todos esses vetores para encontrar a espécie correta?} O BirdSearch
investiga duas respostas distintas, ilustradas na Figura~\ref{fig:fusao}.

\begin{figure}[ht]
\centering
\begin{tikzpicture}[
    seg/.style={rectangle, rounded corners=2pt,
                draw=azulRI!60, fill=azulRI!12,
                minimum width=1.0cm, minimum height=0.6cm,
                font=\tiny\bfseries, text centered},
    emb/.style={rectangle, rounded corners=2pt,
                draw=verdeRI!60, fill=verdeRI!10,
                minimum width=1.0cm, minimum height=0.5cm,
                font=\tiny, text centered},
    idx/.style={cylinder, shape border rotate=90,
                draw=azulRI!50, fill=azulRI!5,
                minimum width=1.1cm, minimum height=0.7cm,
                font=\tiny\bfseries, text centered, aspect=0.25},
    res/.style={rectangle, rounded corners=2pt,
                draw=orange!60, fill=orange!10,
                minimum width=2.2cm, minimum height=0.5cm,
                font=\tiny, text centered},
    sup/.style={rectangle, rounded corners=2pt,
                draw=purple!60, fill=purple!10,
                minimum width=1.1cm, minimum height=0.6cm,
                font=\tiny\bfseries, text centered},
    arr/.style={-Stealth, thick, gray!60},
    lbl/.style={font=\scriptsize\bfseries},
    note/.style={font=\tiny\itshape, color=cinzaRI, align=center}
  ]

  % ── E1: Late Fusion (top half) ──────────────────────────────────────────
  \node[lbl, color=azulRI] at (-0.2, 5.8) {E1 --- Fusão Tardia};

  % Audio bar
  \draw[thick, azulRI!40, fill=azulRI!5, rounded corners=2pt]
    (0.0, 5.1) rectangle (6.2, 5.5);
  \node[font=\tiny\color{cinzaRI}] at (3.1, 5.3) {áudio de consulta (ex: 15 s)};

  % Segments with overlap
  \foreach \x/\n in {0.1/S1, 1.35/S2, 2.6/S3, 3.85/S4} {
    \node[seg] (s\n) at (\x+0.5, 4.6) {\n};
  }
  % overlap braces
  \draw[decorate, decoration={brace,amplitude=3pt}, gray!50]
    (1.25,4.35) -- (0.75,4.35) node[midway,below,font=\tiny\color{cinzaRI}]{hop};

  % Embeddings
  \foreach \x/\n in {0.1/S1, 1.35/S2, 2.6/S3, 3.85/S4} {
    \node[emb] (e\n) at (\x+0.5, 3.7) {$\mathbf{e}_\n$};
    \draw[arr] (s\n) -- (e\n);
  }

  % Index (shared corpus)
  \node[idx] (idxE1) at (7.2, 4.15) {Índice\\HNSW};

  % 4 searches
  \foreach \n/\col in {S1/azulRI,S2/verdeRI,S3/orange,S4/purple} {
    \draw[-Stealth, \col!60, thick] (e\n.east) -- (idxE1.west)
      node[midway, above, font=\tiny\color{\col}] {};
  }

  % 4 result lists
  \node[res] (r1) at (5.5, 3.05) {Lista 1: \textit{Tangara}, \ldots};
  \node[res] (r2) at (5.5, 2.55) {Lista 2: \textit{Tangara}, \ldots};
  \node[res] (r3) at (5.5, 2.05) {Lista 3: \textit{Pipra}, \ldots};
  \node[res] (r4) at (5.5, 1.55) {Lista 4: \textit{Tangara}, \ldots};
  \foreach \n in {1,2,3,4} {
    \draw[arr] (idxE1.south) |- (r\n.east);
  }

  % Combiner
  \node[sup, fill=orange!15, draw=orange!60, minimum width=1.5cm]
    (comb) at (3.5, 1.55) {Reranker};
  \foreach \n in {1,2,3,4} { \draw[arr] (r\n.west) -- (comb.east); }

  \node[res, fill=verdeRI!15, draw=verdeRI!60] (out1) at (1.5, 1.55)
    {Ranking Final};
  \draw[arr] (comb.west) -- (out1.east);

  \node[note] at (3.6, 0.95)
    {4 buscas independentes $\rightarrow$ combina 4 listas de resultados};

  % ── Separator ───────────────────────────────────────────────────────────
  \draw[dashed, gray!40] (-0.3, 0.6) -- (8.2, 0.6);

  % ── E2: Early Fusion (bottom half) ──────────────────────────────────────
  \node[lbl, color=verdeRI] at (-0.2, 0.2) {E2 --- Fusão Precoce};

  % Audio bar
  \draw[thick, verdeRI!40, fill=verdeRI!5, rounded corners=2pt]
    (0.0, -0.5) rectangle (6.2, -0.1);
  \node[font=\tiny\color{cinzaRI}] at (3.1, -0.3) {mesmo áudio de consulta (15 s)};

  % Segments
  \foreach \x/\n in {0.1/S1, 1.35/S2, 2.6/S3, 3.85/S4} {
    \node[seg] (q\n) at (\x+0.5, -1.0) {\n};
  }

  % Embeddings
  \foreach \x/\n in {0.1/S1, 1.35/S2, 2.6/S3, 3.85/S4} {
    \node[emb] (qe\n) at (\x+0.5, -1.75) {$\mathbf{e}_\n$};
    \draw[arr] (q\n) -- (qe\n);
  }

  % Max-pool aggregation
  \node[sup] (pool) at (2.85, -2.55) {max-pool $\rightarrow$ $\mathbf{q}_{\text{super}}$};
  \foreach \n in {S1,S2,S3,S4} { \draw[arr] (qe\n.south) -- (pool.north); }

  % Index
  \node[idx] (idxE2) at (5.2, -2.55) {Índice\\Plano};
  \draw[arr] (pool.east) -- (idxE2.west)
    node[midway, above, font=\tiny] {1 busca};

  % Single result
  \node[res, fill=verdeRI!15, draw=verdeRI!60] (out2) at (7.2, -2.55)
    {Ranking Final};
  \draw[arr] (idxE2.east) -- (out2.west);

  \node[note] at (3.6, -3.3)
    {segmentos fundidos \textbf{antes} da busca $\rightarrow$ 1 única busca no índice};

\end{tikzpicture}
\caption{Comparação entre as duas estratégias de fusão. Em E1 (fusão tardia),
cada segmento da consulta faz uma busca independente no índice e as listas
resultantes são combinadas pelo reranker. Em E2 (fusão precoce), todos os
embeddings são fundidos em um único super-embedding por max-pooling antes
de uma única busca.}
\label{fig:fusao}
\end{figure}

\subsubsection{Estratégia 1 --- Fusão Tardia (E1)}

Na fusão tardia, o sistema trata \textbf{cada segmento da consulta como uma
consulta independente}. Para um áudio de 15~s com sobreposição de 50\%, isso
gera 4 segmentos de 5~s e, portanto, 4 buscas no índice HNSW. Cada busca
retorna uma lista dos $k$ segmentos do corpus mais similares àquele trecho
específico.

O raciocínio por trás dessa abordagem é direto: se o pássaro cantou de forma
clara apenas nos primeiros 5~s do áudio, o segmento S1 já é suficiente para
identificá-lo. Os demais segmentos podem conter silêncio, ruído ou outra
espécie, mas suas listas de resultados também são consideradas. O reranker
(Seção~\ref{sub:reranking}) então combina todas as listas em um ranking final,
dando mais peso às evidências que aparecem consistentemente em múltiplos
segmentos.

\textbf{Vantagem principal}: preserva a diversidade de evidências. Se o pássaro
cantou em diferentes momentos do áudio, cada segmento captura um aspecto
diferente, e a combinação das listas tende a cobrir mais candidatos relevantes
(daí o P@5 maior).

\textbf{Custo}: $N_q$ buscas no índice ($N_q \approx 4$--$12$), o que eleva
a latência para 22,6~ms na média. Com 363.917 documentos no índice HNSW, cada
busca individualmente é rápida ($O(\log N)$), mas o custo se multiplica.

\subsubsection{Estratégia 2 --- Fusão Precoce (E2)}

Na fusão precoce, os embeddings de todos os segmentos da consulta são
\textbf{fundidos em um único vetor antes de qualquer busca}. A operação de
fusão usada é o \textit{max-pooling}: para cada uma das 1280 dimensões do
embedding, é mantido o valor máximo encontrado entre todos os segmentos:
\begin{equation}
  \mathbf{q}_{\text{super}}[d] = \max_{j \in \{1,\ldots,N_q\}}
  \mathbf{e}_j[d], \quad d = 1,\ldots,1280
  \label{eq:super}
\end{equation}
O vetor resultante é então normalizado em L2. Intuitivamente, o max-pooling
preserva, em cada dimensão do espaço de embedding, a \emph{evidência mais
forte} encontrada em qualquer trecho do áudio. Esse vetor comprimido representa
o ``melhor de cada segmento'' e é usado como consulta única no índice plano.

\textbf{Vantagem principal}: uma única busca no índice reduz a latência para
6,3~ms na média (3,5$\times$ mais rápido que E1). O índice também é menor:
27.309 vetores (um por arquivo de treino) contra 363.917 (um por segmento).

\textbf{Custo}: ao fundir tudo em um vetor, informações sobre \emph{quando}
a vocalização ocorreu são perdidas. Se dois segmentos correspondem a espécies
diferentes no mesmo áudio, o super-embedding pode misturar seus sinais.

\subsection{Reranking}
\label{sub:reranking}

Após a busca kNN retornar os candidatos mais próximos, o reranking é a etapa
que \textbf{reordena essa lista} para colocar a espécie correta o mais próximo
do topo possível. É análogo ao que um motor de busca faz ao ordenar páginas
da web: a recuperação vetorial encontra candidatos plausíveis rapidamente;
o reranker decide qual deles é mais relevante com critérios mais finos.

No BirdSearch, o reranking é especialmente importante em E1 (fusão tardia),
porque o sistema recebe \emph{múltiplas listas de candidatos} --- uma por
segmento --- e precisa combiná-las em uma única lista ordenada.
A Figura~\ref{fig:reranking} ilustra como esse processo funciona.

\begin{figure}[ht]
\centering
\begin{tikzpicture}[
    seg/.style={rectangle, rounded corners=2pt,
                draw=azulRI!50, fill=azulRI!8,
                minimum width=0.9cm, minimum height=0.55cm,
                font=\tiny\bfseries, text centered},
    lst/.style={rectangle, rounded corners=2pt,
                draw=gray!50, fill=gray!6,
                minimum width=2.5cm, minimum height=0.45cm,
                font=\tiny, text centered, inner sep=2pt},
    rk/.style={rectangle, rounded corners=3pt,
               draw=orange!60, fill=orange!10,
               minimum width=2.2cm, minimum height=0.6cm,
               font=\scriptsize\bfseries, text centered},
    out/.style={rectangle, rounded corners=3pt,
                draw=verdeRI!60, fill=verdeRI!10,
                minimum width=2.8cm, minimum height=0.45cm,
                font=\tiny, text centered},
    arr/.style={-Stealth, thick, black!55},
    note/.style={font=\tiny\itshape, color=cinzaRI}
  ]

  % Segments column
  \node[seg] (s1) at (0, 3.8) {S1};
  \node[seg] (s2) at (0, 2.8) {S2};
  \node[seg] (s3) at (0, 1.8) {S3};
  
  \node[note] at (0, 0.9) {segmentos da consulta};

  % kNN arrow

  % Lists
  \node[lst, text=azulRI!80] (l11) at (2.8, 4.25) {1. \textit{Tangara} 0,97};
  \node[lst] (l12) at (2.8, 3.75) {2. \textit{Dacnis} 0,91};
  \node[lst] (l13) at (2.8, 3.25) {3. \textit{Pipra} 0,88};
  \node[note, align=center] at (2.8, 2.9) {Lista S1};

  \node[lst, text=azulRI!80] (l21) at (2.8, 2.45) {1. \textit{Tangara} 0,95};
  \node[lst] (l22) at (2.8, 1.95) {2. \textit{Tangara} 0,89};
  \node[lst] (l23) at (2.8, 1.45) {3. \textit{Chlorophanes} 0,82};
  \node[note, align=center] at (2.8, 1.1) {Lista S2};

  \node[lst, text=azulRI!80] (l31) at (2.8, 0.65) {1. \textit{Tangara} 0,93};
  \node[lst] (l32) at (2.8, 0.15) {2. \textit{Dacnis} 0,87};
  \node[lst, text=azulRI!80] (l33) at (2.8,-0.35) {3. \textit{Tangara} 0,85};
  \node[note, align=center] at (2.8,-0.75) {Lista S3};

  % Arrows from segments to lists
  
  \draw[arr] (s1.east) -- (l11.west);
  \draw[arr] (s2.east) -- (l21.west);
  \draw[arr] (s3.east) -- (l31.west);

  % Reranker box
  \node[rk] (rkbox) at (5.5, 1.9) {Reranker};

  % Arrows from lists to reranker
  \foreach \lst in {l11,l12,l13,l21,l22,l23,l31,l32,l33} {
    \draw[arr] (\lst.east) -- (rkbox.west);
  }

  % Output ranking
  \node[out, fill=verdeRI!20, draw=verdeRI!70, font=\small\bfseries]
    (o1) at (9.0, 3.3) {1\textordmasculine\ \textit{Tangara} spp.};
  \node[out] (o2) at (8.9, 2.7) {2\textordmasculine\ \textit{Dacnis} spp.};
  \node[out] (o3) at (8.8, 2.1) {3\textordmasculine\ \textit{Pipra} spp.};
  \node[out] (o4) at (9.1, 1.5) {4\textordmasculine\ \textit{Chlorophanes}};
  \node[note] at (9.0, 0.9) {Ranking final};

  \draw[arr] (rkbox.east) -- (o1.west);
  \draw[arr] (rkbox.east) -- (o2.west);
  \draw[arr] (rkbox.east) -- (o3.west);
  \draw[arr] (rkbox.east) -- (o4.west);


\end{tikzpicture}
\caption{Processo de reranking em E1. Três segmentos da consulta geram três
listas de candidatos via busca kNN. O reranker consolida todas as listas em
um único ranking final, usando diferentes estratégias para combinar as evidências.
Neste exemplo, \textit{Tangara} aparece consistentemente nas três listas e
sobe ao topo.}
\label{fig:reranking}
\end{figure}

Os 14 rerankers implementados organizam-se em cinco famílias, da mais simples
à mais sofisticada:

\textbf{(1) Agregação por score} (\texttt{max}, \texttt{mean},
\texttt{topk\_mean}, \texttt{weighted\_topk}, \texttt{median}, \texttt{threshold}):
olham diretamente os valores de similaridade. O \texttt{max}, por exemplo,
atribui a cada espécie candidata o maior score encontrado em qualquer segmento.
É simples e eficaz quando um único segmento já contém a vocalização clara.
O \texttt{topk\_mean} usa a média dos $k=3$ maiores scores, equilibrando
evidências fortes com consistência.

\textbf{(2) Votação} (\texttt{hit}): ignora completamente os valores dos
scores e conta apenas \emph{quantas vezes} cada espécie apareceu entre os
candidatos das listas. Uma espécie que aparece em 3 das 4 listas recebe
pontuação 3. Funciona bem quando a concordância entre segmentos é um sinal
mais confiável que os scores absolutos.

\textbf{(3) Fusão por rank} (\texttt{rrf}, \texttt{borda}): usam apenas a
\emph{posição} das espécies nas listas, sem considerar os valores dos scores
(Eq.~\eqref{eq:rrf}). O RRF calcula $1/(60 + \text{posição})$ para cada
aparição de uma espécie e soma essas contribuições. O valor $k=60$ é uma
constante empírica que suaviza a diferença entre posições altas. O Borda Count
faz algo similar: dá mais pontos a quem aparece em posições melhores.

\textbf{(4) Score normalizado} (\texttt{softmax}, \texttt{attention}): antes
de combinar as listas, normalizam os scores de cada lista separadamente via
softmax com temperatura $\tau$. Com $\tau=0{,}05$--$0{,}1$, o softmax é
muito concentrado: quem está no topo de cada lista recebe quase todo o peso,
enquanto resultados medianos são praticamente ignorados:
\begin{equation}
  s_{\text{att}}(\ell) = \sum_j \hat{w}_j
    \underbrace{\sum_{i:\,\text{label}=\ell}
      \frac{e^{s_{ji}/\tau}}{\sum_{i'} e^{s_{ji'}/\tau}}}_{\text{softmax por lista}},
  \quad
  \hat{w}_j = \frac{e^{\max_i s_{ji}/\tau}}{\sum_{j'} e^{\max_{i} s_{j'i}/\tau}}
  \label{eq:att_rk}
\end{equation}
O \texttt{attention} vai além e também pondera as próprias consultas: segmentos
em que o sistema está mais ``confiante'' (pico de similaridade mais alto)
recebem mais peso do que segmentos ruidosos. Esse mecanismo duplo de atenção
--- sobre os candidatos e sobre as consultas --- explica por que é o melhor
reranker puro nos experimentos.

\textbf{(5) Taxonômico e Híbridos} (\texttt{taxonomy\_boost} e variantes
híbridas): discutidos em detalhes na Seção~\ref{sub:tax} em conjunto com as
métricas taxonômicas. Os híbridos combinam dois ou três rerankers via soma
ponderada de scores normalizados e operam sobre listas já computadas, sem
nova busca no índice, atingindo latência de 0,3--1,0~ms.

\subsection{Reranking Taxonômico e Métricas Biológicas}
\label{sub:tax}

O \texttt{taxonomy\_boost} é um reranker de segunda camada: ele parte de uma
pontuação base (calculada por qualquer outro reranker) e aplica um
\textbf{bônus taxonômico} a espécies biologicamente próximas da espécie
mais bem pontuada $\hat\ell$:
\begin{equation}
  s_{\text{tax}}(\ell) = s_{\text{base}}(\ell)
    + 0{,}20 \cdot \mathbf{1}[\text{genus}(\ell) = \text{genus}(\hat\ell)]
    + 0{,}10 \cdot \mathbf{1}[\text{tokens}(\ell) \cap \text{tokens}(\hat\ell) \neq \emptyset]
  \label{eq:tax}
\end{equation}

Para entender como funciona na prática, considere um exemplo concreto.
Suponha que a busca retornou estes candidatos (com scores base normalizados):

\begin{center}
\begin{tabular}{lccc}
\toprule
\textbf{Espécie candidata} & \textbf{Gênero} & \textbf{Score base} & \textbf{Score após boost} \\
\midrule
\textit{Tangara cyanicollis}  & \textit{Tangara} & 0,80 & 0,80 (topo, $\hat\ell$) \\
\textit{Chlorophanes spiza}   & \textit{Chlorophanes} & 0,78 & 0,78 \\
\textit{Tangara chilensis}    & \textit{Tangara} & 0,75 & \textbf{0,95} (+0,20) \\
\textit{Tangara mexicana}     & \textit{Tangara} & 0,71 & \textbf{0,91} (+0,20) \\
\textit{Blue Dacnis}          & \textit{Dacnis}  & 0,69 & 0,69 \\
\bottomrule
\end{tabular}
\end{center}

As espécies do mesmo gênero (\textit{Tangara}) recebem um bônus de $+0{,}20$
e sobem no ranking, mesmo com score base menor. Isso reflete a expectativa de
que cantos de espécies do mesmo gênero são mais similares entre si do que
cantos de gêneros distantes, tornando esses erros ``mais aceitáveis''.

Os \textbf{rerankers híbridos} combinam essa lógica taxonômica com rerankers
de score em uma soma ponderada. Por exemplo, o
\texttt{hybrid\_attention\_softmax\_taxonomy} calcula:
\[
  s_{\text{hib}}(\ell) = 0{,}45\,\hat{s}_{\text{att}}(\ell)
                       + 0{,}35\,\hat{s}_{\text{sm}}(\ell)
                       + 0{,}20\,\hat{s}_{\text{tax}}(\ell)
\]
onde $\hat{s}$ indica score normalizado para $[0,1]$. A normalização impede
que um ranker com escalas absolutas maiores domine os demais. Esse híbrido
opera \emph{inteiramente sobre as listas já computadas}, sem nova chamada ao
índice, atingindo latência de 0,3--1,0~ms.

\subsection{Dataset e Configuração Experimental}

O corpus é o \textbf{BirdCLEF 2026 Training Set}: 34.144 gravações de campo,
571 espécies, filtradas por qualidade $\geq 3{,}0$ estrelas (escala Xeno-canto).
A divisão treino/teste é estratificada por espécie na proporção 80/20.

\begin{center}
\begin{tabular}{lcc}
\toprule
\textbf{Conjunto} & \textbf{Arquivos} & \textbf{Segmentos (overlap 50\%)} \\
\midrule
Treino & 27.309 & 363.917 (E1) / 27.309 (E2) \\
Teste  &  6.835 & --- (consultas) \\
\bottomrule
\end{tabular}
\end{center}

Quatro experimentos principais foram conduzidos: (1) \textit{strategy\_compare}
--- comparação direta E1 \textit{vs.}\ E2 com Perch v2; (2)
\textit{strategy\_compare\_giant} --- seis variantes com/sem overlap e com/sem
\texttt{spectral\_gating}; (3) \textit{ranking\_suite\_advanced} --- sweep de
14 rerankers sobre E1 e E2 com Perch v2, incluindo métricas taxonômicas; (4)
\textit{ranking\_sweep\_all} --- sweep de rerankers sobre E1 e E2 com BirdNET v3.0.

% ══════════════════════════════════════════════════════════════════════════════
\section{Resultados}
\label{sec:res}
% ══════════════════════════════════════════════════════════════════════════════

\subsection{Comparação Base E1 \textit{vs.} E2}

A Tabela~\ref{tab:base} apresenta a comparação direta entre as estratégias
base (Perch v2, dataset completo, 6.835 consultas).

\begin{table}[H]
\centering
\caption{Comparação base E1 (HNSW + RRF) \textit{vs.}\ E2 (Flat + Max) com
Perch v2. Dataset completo, 6.835 consultas.}
\label{tab:base}
\begin{tabular}{lccc}
\toprule
\textbf{Métrica} & \textbf{E1 --- Segmentos} & \textbf{E2 --- Super-emb} & \textbf{$\Delta$} \\
\midrule
MAP              & 0,8191 & \textbf{0,8227} & +0,0036 \\
MRR              & 0,8191 & \textbf{0,8227} & +0,0036 \\
nDCG             & 0,8417 & \textbf{0,8457} & +0,0040 \\
P@1              & 0,7697 & \textbf{0,7729} & +0,0032 \\
P@5              & \textbf{0,2518} & 0,2207  & --0,0311 \\
Latência (média) & 22,6~ms & \textbf{6,3~ms} & $3{,}6\times$ \\
Latência (p95)   & 33,3~ms & \textbf{6,4~ms} & $5{,}2\times$ \\
\bottomrule
\end{tabular}
\end{table}

A E2 supera E1 em MAP, MRR, nDCG e P@1, ao passo que E1 mantém vantagem
em P@5 (0,2518 \textit{vs.}\ 0,2207): a fusão tardia preserva mais candidatos
relevantes na lista expandida, o que é útil em aplicações que exibem os
5 principais resultados. A redução de latência de E2 é expressiva ---
$3{,}6\times$ na média e $5{,}2\times$ no p95 --- pois E2 emite uma única
busca no índice plano (27.309 vetores) contra $N_q \approx 6$--$8$ buscas
no HNSW de E1 (363.917 vetores).

\subsection{Impacto da Sobreposição e do Pré-processamento}

\begin{table}[H]
\centering
\caption{Impacto da sobreposição e do filtro espectral. Experimento
\textit{strategy\_compare\_giant}, Perch v2, 6.835 consultas.}
\label{tab:giant}
\begin{tabular}{lccccc}
\toprule
\textbf{Variante} & \textbf{Overlap} & \textbf{Noise filter} & \textbf{MAP} & \textbf{P@1} & \textbf{Lat.\ (ms)} \\
\midrule
E1 --- overlap, sem noise    & Sim & Não & 0,8194 & 0,7697 & 21,9 \\
E2 --- overlap, sem noise    & Sim & Não & \textbf{0,8227} & \textbf{0,7729} & \textbf{6,2} \\
E1 --- sem overlap, sem noise & Não & Não & 0,7980 & 0,7424 & 14,3 \\
E2 --- sem overlap, sem noise & Não & Não & 0,7929 & 0,7397 & 6,3 \\
E1 --- sem overlap, com noise & Não & Sim & 0,7489 & 0,6865 & 14,4 \\
E2 --- sem overlap, com noise & Não & Sim & 0,7607 & 0,7067 & 6,3 \\
\bottomrule
\end{tabular}
\end{table}

A Tabela~\ref{tab:giant} revela três padrões:

\textbf{(1) A sobreposição é essencial.} Removê-la reduz o MAP em 2,1 pontos
em E1 (0,8194\,$\to$\,0,7980) e 3,0 pontos em E2 (0,8227\,$\to$\,0,7929).
Vocalizações que cruzam fronteiras de janela só são representadas integralmente
com overlap.

\textbf{(2) O \texttt{spectral\_gating} degrada o desempenho neste contexto.}
Quando ativado sobre E1 sem overlap, o MAP cai a 0,7489 --- uma queda de 4,9
pontos em relação ao E1 com overlap. O Perch v2, treinado em áudio de campo
bruto, incorporou robustez a ruído em seus pesos; o pré-processamento espectral
remove componentes que o modelo usa como pistas discriminativas.

\textbf{(3) A latência de E2 é estável.} Em todas as variantes, E2 permanece
entre 6,2 e 6,3~ms, enquanto E1 varia entre 14 e 22~ms conforme o número de
segmentos por consulta.

\subsection{Sweep de Rerankers com Perch v2}

\begin{table}[H]
\centering
\caption{Resultados do sweep de rerankers com Perch v2 (6.835 consultas).
Métricas taxonômicas medidas em E1. Melhores por estratégia em \textbf{negrito}.}
\label{tab:sweep_perch}
\begin{tabular}{llccccc}
\toprule
\textbf{Est.} & \textbf{Ranker} & \textbf{MAP} & \textbf{nDCG} & \textbf{P@1} & \textbf{genus@5} & \textbf{Lat.} \\
\midrule
E1 & \texttt{attention} ($\tau=0{,}05$)   & 0,8689 & \textbf{0,8862} & \textbf{0,8269} & 0,9223 & 12,5~ms \\
E1 & \texttt{softmax} ($\tau=0{,}1$)      & 0,8682 & 0,8856 & 0,8252 & 0,9242 & 12,6~ms \\
E1 & \texttt{hybrid\_att\_sm\_tax}         & \textbf{0,8691} & \textbf{0,8862} & 0,8263 & \textbf{0,9247} & \textbf{1,0~ms} \\
E1 & \texttt{taxonomy\_boost}             & 0,8628 & 0,8813 & 0,8180 & 0,9207 & 12,0~ms \\
E1 & \texttt{hybrid\_att\_rrf\_tax}        & 0,8622 & 0,8808 & 0,8202 & 0,9163 & 0,8~ms \\
E1 & \texttt{rrf} ($k=60$)                & 0,8179 & 0,8406 & 0,7691 & 0,8844 & 11,9~ms \\
E1 & \texttt{borda}                        & 0,8170 & 0,8403 & 0,7672 & 0,8832 & 11,9~ms \\
E1 & \texttt{topk\_mean}                   & 0,7632 & 0,7890 & 0,7071 & 0,8375 & 12,0~ms \\
\midrule
E2 & \texttt{attention} ($\tau=0{,}05$)   & \textbf{0,8494} & \textbf{0,8694} & \textbf{0,8041} & 0,9131 & 5,1~ms \\
E2 & \texttt{hybrid\_att\_sm\_tax}         & 0,8467 & 0,8665 & 0,7968 & \textbf{0,9156} & \textbf{0,3~ms} \\
E2 & \texttt{hybrid\_att\_rrf\_tax}        & 0,8476 & 0,8677 & 0,8029 & 0,9069 & \textbf{0,3~ms} \\
E2 & \texttt{softmax}                      & 0,8436 & 0,8656 & 0,7925 & 0,9141 & 5,0~ms \\
E2 & \texttt{taxonomy\_boost}             & 0,8200 & 0,8467 & 0,7561 & 0,9055 & 5,2~ms \\
E2 & \texttt{rrf}                          & 0,8227 & 0,8457 & 0,7729 & 0,8907 & 5,0~ms \\
\bottomrule
\end{tabular}
\end{table}

A Tabela~\ref{tab:sweep_perch} mostra resultados consistentes em ambas as
estratégias. Em E1, \texttt{attention} (MAP\,=\,0,8689) e \texttt{softmax}
(0,8682) superam o \texttt{rrf} base (0,8179) em $\approx$\,5 pontos de MAP,
evidenciando que os scores absolutos do Perch v2 carregam informação além da
posição. O híbrido \texttt{hybrid\_att\_sm\_tax} atinge MAP\,=\,0,8691 com
latência de apenas 1,0~ms --- porque opera sobre as listas pré-computadas do
índice base, sem busca adicional.

Em E2, o \texttt{attention} (0,8494) lidera, com o híbrido
\texttt{hybrid\_att\_rrf\_tax} logo atrás (0,8476) a 0,3~ms. O
\texttt{taxonomy\_boost} tem o melhor genus@5 relativo em E2 (0,9055),
sendo preferível em aplicações onde a relevância taxonômica importa mais
que a precisão exata.

\subsection{Sweep de Rerankers com BirdNET v3.0}

\begin{table}[H]
\centering
\caption{Resultados selecionados com BirdNET v3.0 (6.835 consultas).
363.917 documentos em E1, 27.309 em E2. genus@5 não foi coletado
neste experimento (indicado por ---). Melhores por estratégia em \textbf{negrito}.}
\label{tab:sweep_birdnet}
\begin{tabular}{llccccc}
\toprule
\textbf{Est.} & \textbf{Ranker} & \textbf{MAP} & \textbf{nDCG} & \textbf{P@1} & \textbf{genus@5} & \textbf{Lat.} \\
\midrule
E1 & \texttt{attention}       & \textbf{0,8944} & \textbf{0,9105} & \textbf{0,8549} & --- & 18,7~ms \\
E1 & \texttt{softmax}         & 0,8910 & 0,9078 & 0,8506 & --- & 22,4~ms \\
E1 & \texttt{taxonomy\_boost} & 0,8894 & 0,9064 & 0,8478 & --- & 18,5~ms \\
E1 & \texttt{rrf}             & 0,8312 & 0,8555 & 0,7775 & --- & 20,3~ms \\
E1 & \texttt{topk\_mean}      & 0,7707 & 0,8003 & 0,7083 & --- & 16,7~ms \\
E1 & \texttt{hit}             & 0,6743 & 0,7270 & 0,5649 & --- & 16,8~ms \\
E1 & \texttt{mean}            & 0,2520 & 0,2972 & 0,1671 & --- & 17,3~ms \\
\midrule
E2 & \texttt{attention}       & \textbf{0,8373} & \textbf{0,8604} & \textbf{0,7881} & --- & 6,2~ms \\
E2 & \texttt{softmax}         & 0,8298 & 0,8540 & 0,7786 & --- & 6,0~ms \\
E2 & \texttt{taxonomy\_boost} & 0,8240 & 0,8493 & 0,7700 & --- & 6,2~ms \\
E2 & \texttt{rrf}             & 0,8108 & 0,8371 & 0,7557 & --- & 5,9~ms \\
E2 & \texttt{mean}            & 0,5163 & 0,5973 & 0,3557 & --- & 6,0~ms \\
\bottomrule
\end{tabular}
\end{table}

Com BirdNET, o \texttt{attention} em E1 alcança MAP\,=\,0,8944 --- o
\textbf{melhor resultado de todo o projeto} (Tabela~\ref{tab:geral}),
superando o Perch v2 em 2,55 pontos percentuais. Isso sugere que os embeddings
do BirdNET possuem maior margem entre espécies similares, amplificando os
benefícios de rankers que exploram scores absolutos.

O fenômeno mais notável é o colapso dos rankers \texttt{mean} e \texttt{median}
com BirdNET em E1 (MAP\,$\approx$\,0,25): o BirdNET gera scores heterogêneos
entre segmentos, e a média suprime os picos de sinal relevantes. O
\texttt{attention} e o \texttt{softmax}, que concentram o peso nos picos via
softmax com $\tau=0{,}05$--$0{,}1$, são imunes a esse colapso.

\subsection{Comparação Geral de Desempenho}

A Tabela~\ref{tab:geral} consolida os melhores rerankers de cada modelo e
estratégia de fusão. Para cada linha escolheu-se o reranker que maximiza o
MAP dentre os avaliados na suíte completa.
O genus@5 do BirdNET foi obtido executando
\texttt{experiments/ranking\_suite\_birdnet.py}, que replica a mesma lógica
de avaliação taxonômica da suíte Perch~v2.

\begin{table}[ht]
\centering
\caption{Comparação geral: melhores rerankers por modelo e estratégia.
PV\,=\,Perch v2; BN\,=\,BirdNET v3.0; E1\,=\,fusão tardia (HNSW);
E2\,=\,fusão antecipada (Flat). Lat.\ = latência média de busca (ms).
$\dagger$\,genus@5 requer execução de
\texttt{ranking\_suite\_birdnet.py}.}
\label{tab:geral}
\begin{tabular}{llccccc}
\toprule
\textbf{Modelo} & \textbf{Est.} & \textbf{MAP} & \textbf{nDCG} & \textbf{P@1} & \textbf{genus@5} & \textbf{Lat.\ (ms)} \\
\midrule
\multirow{2}{*}{Perch v2}
  & E1 & \textbf{0,8691} & \textbf{0,9125} & \textbf{0,8311} & \textbf{0,9247} &  1,0 \\
  & E2 & 0,8495          & 0,8906          & 0,8060          & 0,9071          &  0,3 \\
\midrule
\multirow{2}{*}{BirdNET v3.0}
  & E1 & 0,8944          & 0,9105          & 0,8549          & $\dagger$       & 18,7 \\
  & E2 & 0,8373          & 0,8814          & 0,7998          & $\dagger$       &  6,2 \\
\bottomrule
\end{tabular}
\end{table}

\textbf{E1 supera E2 em ambos os modelos.} A fusão tardia, ao emitir uma
consulta por segmento e combinar as listas com RRF ou \texttt{attention},
captura variações intra-gravação que a super-embedding de E2 aplana.
O custo é um índice maior (HNSW sobre 363.917 segmentos \textit{vs.}\
Flat sobre 27.309 arquivos) e maior latência.

\textbf{BirdNET lidera em MAP no E1.} O backbone da Cornell, treinado em
6.000+ espécies, generaliza melhor na métrica de precisão média; Perch~v2,
porém, atinge genus@5\,=\,0,925, indicando que seu espaço de embeddings
organiza espécies mais próximas taxonomicamente — resultado coerente com o
treinamento específico em cantos de aves do Xeno-canto~\cite{xeno:14}.

\subsection{Métricas Taxonômicas}

Um diferencial do BirdSearch é medir não só se a espécie exata foi encontrada,
mas o quão biologicamente próximos estão os erros
(Equações~\eqref{eq:genus}--\eqref{eq:common}, Seção~\ref{sub:tax_fund}).

\begin{table}[ht]
\centering
\caption{Métricas taxonômicas em E1 com Perch v2 (6.835 consultas).
genus@5: fração dos top-5 do mesmo gênero; common@10: fração dos top-10
com nome comum sobreposto.}
\label{tab:tax}
\begin{tabular}{lcccc}
\toprule
\textbf{Reranker} & \textbf{MAP} & \textbf{genus@5} & \textbf{genus@10} & \textbf{common@10} \\
\midrule
\texttt{softmax}                  & 0,8682 & \textbf{0,9242} & 0,9419 & 0,9486 \\
\texttt{attention}                & \textbf{0,8689} & 0,9223 & \textbf{0,9424} & \textbf{0,9492} \\
\texttt{hybrid\_att\_sm\_tax}     & 0,8691 & 0,9247 & 0,9423 & 0,9498 \\
\texttt{taxonomy\_boost}          & 0,8628 & 0,9207 & 0,9393 & 0,9482 \\
\texttt{rrf}                      & 0,8179 & 0,8844 & 0,9154 & 0,9254 \\
\texttt{borda}                    & 0,8170 & 0,8832 & 0,9175 & 0,9271 \\
\bottomrule
\end{tabular}
\end{table}

Os resultados da Tabela~\ref{tab:tax} revelam um padrão interessante.

\textbf{Softmax e attention lideram em genus@5.} Embora o
\texttt{taxonomy\_boost} tenha sido projetado explicitamente para promover
espécies do mesmo gênero, os rerankers de score normalizado alcançam valores
semelhantes (0,924 \textit{vs.}\ 0,921) sem usar nenhuma informação taxonômica.
Isso ocorre porque os embeddings do Perch v2, treinados sobre cantos de aves,
aprenderam naturalmente que espécies do mesmo gênero têm cantos similares ---
a taxonomia está ``embutida'' no espaço de embedding.

\textbf{RRF e Borda ficam 3--4 pontos abaixo.} Os rankers de fusão por rank
descartam os valores de similaridade e usam apenas as posições. Ao fazer isso,
perdem a correlação entre similaridade espectral e proximidade taxonômica,
resultando em genus@5\,$\approx$\,0,884. Ou seja, ao não usar os scores, esses
rankers fazem mais erros que não são nem biologicamente próximos.

\textbf{common@10 acima de 0,94 para os melhores rankers.} Isso significa que
em 94\% dos 10 primeiros resultados, o nome comum da espécie tem ao menos uma
palavra em comum com a consulta --- uma indicação de que o sistema raramente
sugere espécies completamente desconexas. Para um ornitólogo em campo, essa
lista de sugestões tem valor mesmo quando a espécie exata não é a primeira.

\subsection{Análise de Latência}

\begin{figure}[H]
\centering
\begin{tikzpicture}
\begin{axis}[
    ybar,
    width=\columnwidth, height=5.5cm,
    bar width=0.55cm,
    ylabel={Latência (ms)},
    xtick=data,
    xticklabels={rrf, attention, softmax, tax\_boost, híbrido\_att\_rrf, híbrido\_att\_sm},
    xticklabel style={rotate=28, anchor=east, font=\scriptsize},
    ymin=0, ymax=28,
    legend style={at={(0.97,0.97)}, anchor=north east, font=\scriptsize},
    grid=major, grid style={dashed, gray!30},
    nodes near coords, nodes near coords align={vertical},
    every node near coord/.append style={font=\tiny},
  ]
  \addplot[fill=azulRI!50, draw=azulRI!80] coordinates {
    (1,11.9)(2,12.5)(3,12.6)(4,12.0)(5,0.8)(6,1.0)
  };
  \addplot[fill=verdeRI!50, draw=verdeRI!80] coordinates {
    (1,5.0)(2,5.1)(3,5.0)(4,5.2)(5,0.3)(6,0.3)
  };
  \legend{E1, E2}
\end{axis}
\end{tikzpicture}
\caption{Latência média de busca por reranker. Híbridos atingem
$<1$~ms pois operam sobre resultados pré-computados. E2 é $\approx$\,2,4$\times$
mais rápido que E1 para rankers simples.}
\label{fig:latency}
\end{figure}

A Figura~\ref{fig:latency} mostra que a latência de rankers simples é
dominada pelo custo da busca no índice ($\approx$\,12~ms em E1, 5~ms em E2).
Rankers híbridos operam sobre as listas já computadas e atingem
0,3--1,0~ms, sendo ideais para implantação com restrições de tempo real.

% ══════════════════════════════════════════════════════════════════════════════
\section{Conclusão}
\label{sec:conc}
% ══════════════════════════════════════════════════════════════════════════════

FAZER AINDA

Como trabalhos futuros, identificamos: (i) investigar filtros de ruído que
preservem as pistas espectrais discriminativas do Perch v2; (ii) explorar
\textit{fine-tuning} dos backbones sobre espécies amazônicas específicas; (iii)
avaliar as representações \texttt{cluster} e \texttt{prototype} para capturar
repertórios vocais multi-modais; e (iv) implementar indexação dinâmica para
suporte a corpus de monitoramento contínuo.

\bibliographystyle{sbc}
\bibliography{sbc-template}

\end{document}
